\thispagestyle{plain}
\chapter*{Resumo}

O desenvolvimento de ficção interativa e jogos baseados em escolhas enfrenta sérios desafios de usabilidade e integridade estrutural à medida que a ramificação narrativa cresce. Embora existam ferramentas populares de autoria como o Twine, a maioria carece de suporte nativo integrado para gestão de localização multi-idioma e de motores de validação estática de caminhos lógicos para detetar inconsistências (como becos sem saída e cenas órfãs) antes da fase de compilação ou execução. O presente trabalho de projeto propõe o \textit{Plot in a Pot} (PiaP), um editor visual e simulador baseado em grafos que unifica a autoria de narrativas escritas no formato aberto Twee 3 (SugarCube) com um motor integrado de tradução de textos e validação estática de caminhos.

Desenvolvida em React e React Flow com uma interface de alto contraste adaptada para acessibilidade visual e cromática, a aplicação integra uma matriz de tradução de textos multi-idioma, um motor de validação estática de caminhos baseado em pesquisa em largura (BFS) e um simulador interativo com testes automáticos suportados por procura de novidade (Smart Walker).

A avaliação qualitativa e quantitativa do sistema foi efetuada através de testes com seis utilizadores com diferentes perfis técnicos e de acessibilidade. Os resultados revelaram uma taxa média de conclusão de tarefas de 91,7\% e uma elevada usabilidade e satisfação, validando a eficácia da ferramenta para autores e game designers na prototipagem de narrativas complexas. Como trabalhos futuros, perspetiva-se a execução de IA local via WebGPU, edição colaborativa em tempo real e model checking formal da lógica narrativa.

\keywordspt{Ficção Interativa, Editores de Grafos, Localização de Jogos, Acessibilidade Visual, Validação Estática de Fluxo, Twine, SugarCube.}

\MediaOptionLogicBlank

\pdfbookmark[1]{Abstract}{abstract}
\chapter*{Abstract}

The development of interactive fiction and choice-based games faces serious challenges in usability and structural integrity as the narrative branching expands. Although popular authoring tools like Twine exist, most lack integrated native support for multi-language localization management and static path-validation engines to detect inconsistencies (such as dead ends and orphan scenes) before compilation or execution. This project work proposes \textit{Plot in a Pot} (PiaP), a graph-based visual editor and simulator that unifies the authoring of interactive stories written in the open Twee 3 format (SugarCube) with an integrated translation matrix and static path-validation engine.

Developed in React and React Flow with a high-contrast interface tailored for visual and chromatic accessibility, the application integrates a multi-language translation matrix, a static path-validation engine based on Breadth-First Search (BFS), and an interactive simulator with automated testing powered by novelty search (Smart Walker).

Qualitative and quantitative evaluation was conducted through user testing with six participants of varying technical backgrounds and accessibility needs. The results showed a 91.7\% average task completion rate alongside high levels of usability and satisfaction, validating the effectiveness of the tool for authors and game designers during narrative prototyping. Future work directions include local AI execution via WebGPU, real-time collaborative editing, and formal model checking of narrative logic.

\keywordsen{Interactive Fiction, Graph Editors, Game Localization, Visual Accessibility, Static Flow Validation, Twine, SugarCube.}

\MediaOptionLogicBlank